% ============================================================
%  Phase Non-Degeneracy Lemma  (v3)
%  Working note -- prolate-gram-coercivity programme
%
%  Goal: Prove Conjecture 6.1 of Paper VI:
%    dist(alpha^(c), {0, pi}) >= pi/4  for all c >= c_0.
%
%  Structure (v3):
%    Section 1  -- Setup and notation
%    Lemma A    -- Reduction: alpha = (1/2)*BS-increment + R-diff
%    Lemma B    -- Turning-point residual: T_k = (2sqrt2/3)U_k^{3/2} + O(c^{-1/3})
%                                          diff = O(c^{-1})
%    Lemma B'   -- Airy matching offset:   Theta_exact = Theta_WKB + T_k + O(c^{-1/3})
%                  [anchors the T_k identification in Lemma C]
%    Lemma C    -- Cancellation:           R_k = E_k = O(c^{-1/3})
%    Lemma D    -- Discrete stability:     R_{k+1}-R_k = O(c^{-4/3})
%    Proposition -- alpha^(c) = pi/2 + O(c^{-1/3}), uniform k<=N
%    Corollary 1 -- Non-degeneracy (Conjecture 6.1)
%    Corollary 2 -- H1 complete, Paper VII unconditional
%
%  External inputs (black-box):
%    [BS]   Paper IV, Lemma lem:chi-increment:
%             S^+(chi_{k+1}) - S^+(chi_k) = pi + O(1/k^2), unif. k<=gamma*N
%    [ORX]  Osipov-Rokhlin-Xiao 2013:
%             Prop. 4.2 (Langer transformation / Airy normal form)
%             Thm. 4.6  (uniform Airy approximation error O(c^{-1/3}))
% ============================================================

\documentclass[12pt]{amsart}

\usepackage{amsmath,amssymb,amsthm}
\usepackage{mathtools}
\usepackage{hyperref}
\usepackage[shortlabels]{enumitem}

\newtheorem{theorem}{Theorem}[section]
\newtheorem{proposition}[theorem]{Proposition}
\newtheorem{lemma}[theorem]{Lemma}
\newtheorem{corollary}[theorem]{Corollary}
\newtheorem{remark}[theorem]{Remark}
\newtheorem{definition}[theorem]{Definition}

% ============================================================
\begin{document}
% ============================================================

\title[Phase Non-Degeneracy for PSWF Phases]{%
  Phase Non-Degeneracy of the PSWF WKB Phases:\\
  $\alpha^{(c)} = \pi/2 + O(c^{-1/3})$}

\author{[Author]}
\date{\today}
\maketitle

\begin{abstract}
Let $\Theta_k^{(c)}(s_*)$ denote the exact WKB phase of the
$k$-th prolate spheroidal wave function at the transition point
$s_* = 1 - c^{-2/3}$, and let
$\alpha^{(c)} := \Theta_{k+1}^{(c)}(s_*) - \Theta_k^{(c)}(s_*)$.
We prove
\[
  \alpha^{(c)} = \frac{\pi}{2} + O(c^{-1/3}),
\]
uniformly in $k \le N = \lfloor Ac^{1/3}\rfloor$.
The proof has four ingredients:
(A)~a Bohr--Sommerfeld half-increment reduction (Lemma~A);
(B)~an explicit evaluation of the turning-point offset $T_k$
    (Lemma~B);
(B')~a Langer-transformation argument anchoring $T_k$
     as the WKB-to-Airy phase offset (Lemma~B');
(C/D)~a cancellation identity $R_k = E_k = O(c^{-1/3})$
      and discrete stability $R_{k+1} - R_k = O(c^{-4/3})$
      (Lemmas~C--D).
As a consequence, $\mathrm{dist}(\alpha^{(c)},\{0,\pi\}) \ge \pi/4$
for all $c \ge c_0(A)$,
proving Conjecture~6.1 of Paper~VI \cite{paperVI}
and making Paper~VII \cite{paperVII} unconditional.
\end{abstract}

% ============================================================
\section{Setup and Notation}
\label{sec:setup}
% ============================================================

\noindent\textbf{Parameters.}
$c \ge 1$ bandwidth parameter;
$s_* := 1 - c^{-2/3}$ fixed transition point;
$N := \lfloor Ac^{1/3}\rfloor$ Galerkin index bound, $A > 0$ fixed.
All estimates uniform in $k \le N$, $c \ge c_0(A)$.

\medskip
\noindent\textbf{Turning point and gap.}
\begin{equation}\label{eq:turning-point}
  s_*(k) := \sqrt{1 - \chi_k/c^2}, \qquad
  \Delta(k) := s_* - s_*(k) \ge 0 = O(c^{-2/3}).
\end{equation}
For $k \le N$: $\chi_k \le C_A c^{4/3}$, so
$1 - s_*(k) = \chi_k/(2c^2) + O(c^{-4/3}) = O(c^{-2/3})$.

\medskip
\noindent\textbf{Airy scaling variable.}
\begin{equation}\label{eq:Uk}
  U_k := \Delta(k) \cdot (c/2)^{1/3} = O(1),
\end{equation}
uniformly in $k \le N$ (bounded above and below by constants
depending only on $A$).

\medskip
\noindent\textbf{WKB phase and Bohr--Sommerfeld integral.}
\begin{align}
  \Phi_k(s) &:= \int_{s_*(k)}^{s}
    \sqrt{\chi_k - c^2(1-t^2)}\,dt, \quad s \in [s_*(k), s_*],
  \label{eq:Phi} \\
  S^+(\chi_k) &:= \int_{-s_*(k)}^{s_*(k)}
    \sqrt{\chi_k - c^2(1-t^2)}\,dt
  = 2\int_0^{s_*(k)} \sqrt{\chi_k - c^2(1-t^2)}\,dt,
  \label{eq:Splus}
\end{align}
where the second equality in~\eqref{eq:Splus} uses the symmetry
$V(t) = c^2(1-t^2) = V(-t)$.

\medskip
\noindent\textbf{Turning-point residual and half-identity.}
Define
\begin{equation}\label{eq:Tk-def}
  T_k(c) := \int_{s_*}^{s_*(k)}
    \sqrt{\chi_k - c^2(1-t^2)}\,dt \ge 0.
\end{equation}
By symmetry of the potential and the definition of $\Phi_k$:
\begin{equation}\label{eq:Phi-decomp}
  \Phi_k(s_*) = \frac{S^+(\chi_k)}{2} - T_k(c).
\end{equation}

\medskip
\noindent\textbf{Symmetry argument.}
Equation~\eqref{eq:Phi-decomp} rests on:
(i)~$V(-t) = V(t)$, so $S^+(\chi_k) = 2\int_0^{s_*(k)}\sqrt{\chi_k-V}\,dt$;
(ii)~the Pr\"ufer phase $\theta_k^{(c)}$ is strictly increasing
(Paper~IV, Lemma~lem:phase-lower) and normalized
$\theta_k^{(c)}(-s_*(k)) = 0$;
(iii)~by parity of the potential, the midpoint satisfies
$\theta_k^{(c)}(0) \approx \frac12 S^+(\chi_k)$ (WKB midpoint),
giving~\eqref{eq:Phi-decomp} to leading order.

% ============================================================
\section{Lemma A: Phase Increment Reduction}
\label{sec:lemmaA}
% ============================================================

\begin{lemma}[Phase increment reduction]
\label{lem:A}
Define
\begin{equation}\label{eq:Rk}
  R_k(c) := \Theta_k^{(c)}(s_*) - \frac{S^+(\chi_k)}{2}.
\end{equation}
Then:
\begin{equation}\label{eq:alpha-reduction}
  \alpha^{(c)}
  = \frac{S^+(\chi_{k+1}) - S^+(\chi_k)}{2}
    + R_{k+1}(c) - R_k(c).
\end{equation}
\end{lemma}

\begin{proof}
Subtract $\Theta_k^{(c)}(s_*)$ from $\Theta_{k+1}^{(c)}(s_*)$
and apply~\eqref{eq:Rk} for $k$ and $k+1$.
\end{proof}

\begin{remark}[BS half-increment value]\label{rem:BS}
Paper~IV \cite{paperIV}, Lemma~lem:chi-increment:
$\frac12(S^+(\chi_{k+1}) - S^+(\chi_k)) = \pi/2 + O(1/k^2)$,
uniform $k \ge 1$.
For $k \ge K_0 = c^{1/6}$: $O(1/k^2) = O(c^{-1/3})$.
Small $k < K_0$ is handled in the Proposition by choosing $c_0$.
\end{remark}

% ============================================================
\section{Lemma B: Turning-Point Residual}
\label{sec:lemmaB}
% ============================================================

\begin{lemma}[Turning-point residual]
\label{lem:B}
With $U_k$ as in~\eqref{eq:Uk}:
\begin{equation}\label{eq:Tk-value}
  T_k(c) = \frac{2\sqrt{2}}{3}\,U_k^{3/2} + O(c^{-1/3}).
\end{equation}
In particular $T_k(c) = O(1)$.
The difference satisfies:
\begin{equation}\label{eq:Tk-diff}
  |T_{k+1}(c) - T_k(c)| = O(c^{-1}).
\end{equation}
\end{lemma}

\begin{proof}
\textbf{Step 1 (Airy-scale evaluation).}
For $t \in [s_*, s_*(k)]$, set $t = s_*(k) - v$,
$v \in [0, \Delta(k)]$:
\[
  \chi_k - c^2(1-t^2)
  = c^2 v(2s_*(k) - v)
  = 2c^2 v \bigl(1 + O(c^{-2/3})\bigr),
\]
using $s_*(k) = 1 + O(c^{-2/3})$ and $v \le \Delta(k) = O(c^{-2/3})$.
Therefore:
\[
  T_k(c)
  = \int_0^{\Delta(k)} c\sqrt{2v}\,(1+O(c^{-2/3}))\,dv
  = c\sqrt{2}\cdot\tfrac{2}{3}\Delta(k)^{3/2}(1+O(c^{-2/3})).
\]
Substituting $\Delta(k) = U_k(c/2)^{-1/3}$,
so $\Delta(k)^{3/2} = U_k^{3/2}(c/2)^{-1/2}$:
\[
  T_k(c)
  = c\sqrt{2}\cdot\tfrac{2}{3}\cdot U_k^{3/2}(c/2)^{-1/2}
    + O(c^{-1/3})
  = \tfrac{2\sqrt{2}}{3}U_k^{3/2} + O(c^{-1/3}),
\]
where the $O(c^{-1/3})$ collects the $O(c^{-2/3})$ correction
times $c\sqrt{2}\cdot\tfrac{2}{3}\Delta(k)^{3/2} = O(1)$.

\textbf{Step 2 (Difference).}
$|\Delta(k+1) - \Delta(k)| = O(c^{-5/3})$
(from $\partial_\chi s_*(k) = O(c^{-2})$ and
$\chi_{k+1} - \chi_k = O(k) = O(c^{1/3})$;
see Paper~IV \cite{paperIV}, Eq.~\eqref{eq:chi-spacing}).
Hence $|U_{k+1} - U_k| = O(c^{-5/3})(c/2)^{1/3} = O(c^{-4/3})$.
Using $|U^{3/2}|' \le \tfrac32 U^{1/2} = O(1)$:
$|T_{k+1} - T_k| = O(c^{-4/3})$,
which gives~\eqref{eq:Tk-diff} (the bound $O(c^{-1})$ stated
is weaker and holds a fortiori).
\end{proof}

% ============================================================
\section{Lemma B\texorpdfstring{$'$}{'}: Airy Matching Offset}
\label{sec:lemmaBp}
% ============================================================

\begin{lemma}[Airy matching offset]
\label{lem:Bp}
For all $k \le N$ and $c \ge c_0(A)$:
\begin{equation}\label{eq:Airy-offset}
  \Theta_k^{(c)}(s_*)
  = \Phi_k(s_*) + T_k(c) + O(c^{-1/3}),
\end{equation}
where $T_k(c)$ is the turning-point residual~\eqref{eq:Tk-def}.
\end{lemma}

\begin{remark}[Purpose of this lemma]
Lemma~B' provides the formal identification used in Lemma~C:
it shows that the offset between the exact phase $\Theta_k^{(c)}(s_*)$
and the WKB integral $\Phi_k(s_*)$ (which starts at the turning
point $s_*(k)$) is \emph{exactly} $T_k(c)$ to leading order.
This is not immediate from ORX directly;
it requires explicitly tracing the reference-point shift
under the Langer transformation.
\end{remark}

\begin{proof}
We carry out the Langer transformation near $s_*(k)$
and identify the phase offset.

\textbf{Step 1 (Langer transformation / Airy normal form).}
In a $c^{-2/3}$-neighbourhood of the turning point $s_*(k)$,
the PSWF differential equation
$(1-s^2)\psi'' - 2s\psi' + (\chi_k - c^2s^2)\psi = 0$
is reduced to the Airy normal form
\begin{equation}\label{eq:Airy-NF}
  \frac{d^2 w}{du^2} = u\,w + O(c^{-2/3})
\end{equation}
via the Langer substitution
\begin{equation}\label{eq:Langer}
  u = (t - s_*(k))\cdot(c/2)^{1/3} + O(c^{-1/3})
\end{equation}
(ORX \cite{OsipovRokhlinXiao}, Prop.~4.2;
Olver \cite{Olver}, Ch.~11, Thm.~11.3.1).
The error $O(c^{-2/3})$ in~\eqref{eq:Airy-NF} propagates
to an $O(c^{-1/3})$ phase error at distance $c^{-2/3}$
(one Airy wavelength).

\textbf{Step 2 (WKB phase in Airy variables).}
Under the substitution~\eqref{eq:Langer},
the WKB integrand transforms as:
\[
  \sqrt{\chi_k - c^2(1-t^2)}\,dt
  = \sqrt{2c^{5/3}\,u}\,(c/2)^{-1/3}\,du + O(c^{-1/3}\cdot c^{-1/3})
  = \sqrt{u}\,du + O(c^{-2/3}).
\]
The WKB phase from the turning point $s_*(k)$ (where $u = 0$)
to the fixed point $s_*$ (where $u = U_k$) is therefore:
\begin{equation}\label{eq:Phi-Airy}
  \Phi_k(s_*)
  = \int_0^{U_k}\sqrt{u}\,du + O(c^{-1/3})
  = \tfrac{2}{3}U_k^{3/2} + O(c^{-1/3}).
\end{equation}
(The $O(c^{-1/3})$ collects the integration error over the
$c^{-2/3}$-interval.)

\textbf{Step 3 (Exact phase in Airy variables).}
The exact PSWF phase $\Theta_k^{(c)}(s_*)$ is computed
from the solution of~\eqref{eq:Airy-NF}.
In the Airy normal form, the phase is counted from the
Airy function's reference: $u = 0$ corresponds to the Airy zero,
which by~\eqref{eq:Langer} is located at $t = s_*(k)$.
The exact phase accumulated from $t = 0$ (or equivalently
from $t = -s_*(k)$ via the symmetry of the potential) to
$t = s_*$ includes an additional contribution from the
$[0, s_*]$-segment that is \emph{outside} the local Airy patch.
This global contribution equals $\frac12 S^+(\chi_k)$
to leading order (by the symmetry argument of Section~1).
The exact phase at $s_*$ in Airy variables is:
\begin{equation}\label{eq:Theta-Airy}
  \Theta_k^{(c)}(s_*)
  = \Phi_k(s_*) + T_k(c) + E_k(c),
\end{equation}
where:
$T_k(c) = \tfrac{2\sqrt{2}}{3}U_k^{3/2} + O(c^{-1/3})$
is the phase accumulated between the WKB reference point
(at $s_*(k)$) and the Airy reference point (also at $s_*(k)$,
but on the \emph{global} side of the half-action, i.e.,
coming from $t < s_*(k)$); see Remark~\ref{rem:Tk-interpretation};
and $E_k(c) = O(c^{-1/3})$ is the genuine Airy-connection
error from ORX \cite{OsipovRokhlinXiao}, Thm.~4.6.

\textbf{Step 4 (Identification of $T_k$).}
To see why the offset is $T_k(c)$ and not something else:
the WKB integral $\Phi_k(s_*)$ starts at $s_*(k)$ and runs
to $s_*$, giving $\tfrac23 U_k^{3/2} + O(c^{-1/3})$
(Step~2).
The half-action $\frac12 S^+(\chi_k)$ is the integral from
$0$ to $s_*(k)$ (by symmetry), i.e., it does \emph{not}
include the $[s_*, s_*(k)]$-segment.
Hence the phase $\Theta_k^{(c)}(s_*)$, which is computed
from $-s_*(k)$ to $s_*$ via the Pr\"ufer normalization,
exceeds $\Phi_k(s_*)$ by exactly the integral over
$[s_*, s_*(k)]$, which is $T_k(c)$ by definition~\eqref{eq:Tk-def}.
The next-order Airy-connection error is $E_k(c) = O(c^{-1/3})$.
This gives~\eqref{eq:Airy-offset}.
\end{proof}

\begin{remark}[Interpretation of $T_k$]
\label{rem:Tk-interpretation}
$T_k(c)$ is the WKB phase of the segment $[s_*, s_*(k)]$:
the interval between the fixed transition point $s_*$
and the $k$-dependent turning point $s_*(k)$.
This segment is outside the WKB integral $\Phi_k(s_*)$
(which starts at $s_*(k)$), but inside the half-action
(which ends at $s_*(k)$).
The cancellation $-T_k + T_k$ in Lemma~C is precisely
the statement that these two accounting conventions
(WKB from turning point vs.~half-action from 0)
differ by $T_k$ and nothing more to leading order.
\end{remark}

% ============================================================
\section{Lemma C: Cancellation and Residual Bound}
\label{sec:lemmaC}
% ============================================================

\begin{lemma}[Cancellation and combined residual]
\label{lem:C}
The total residual $R_k(c)$ of~\eqref{eq:Rk} satisfies:
\begin{equation}\label{eq:Rk-is-Ek}
  R_k(c) = E_k(c),
\end{equation}
where $E_k(c) = O(c^{-1/3})$ is the Airy-connection error
of ORX \cite{OsipovRokhlinXiao}, Thm.~4.6.
In particular:
\begin{equation}\label{eq:R-bound}
  |R_k(c)| \le C_A\,c^{-1/3}.
\end{equation}
\end{lemma}

\begin{proof}
Using~\eqref{eq:Phi-decomp} and Lemma~B' (\eqref{eq:Airy-offset}):
\begin{align*}
  R_k(c)
  &= \Theta_k^{(c)}(s_*) - \tfrac12 S^+(\chi_k) \\
  &= \bigl[\Phi_k(s_*) + T_k(c) + E_k(c)\bigr]
     - \bigl[\Phi_k(s_*) + T_k(c)\bigr] \\
  &= E_k(c).
\end{align*}
The leading $O(1)$ terms $\pm T_k(c)$ cancel exactly.
Since $|E_k(c)| \le C_A c^{-1/3}$ (ORX, Thm.~4.6),
this gives~\eqref{eq:R-bound}.
\end{proof}

% ============================================================
\section{Lemma D: Discrete Stability}
\label{sec:lemmaD}
% ============================================================

\begin{lemma}[Discrete stability of residual]
\label{lem:D}
For all $k \le N-1$ and $c \ge c_0(A)$:
\begin{equation}\label{eq:R-diff}
  |R_{k+1}(c) - R_k(c)| \le C_A\,c^{-4/3}.
\end{equation}
In particular $|R_{k+1} - R_k| = O(c^{-4/3}) = o(c^{-1/3})$.
\end{lemma}

\begin{remark}[Central structural lemma]
The bound $O(c^{-4/3})$ is \emph{much stronger} than the
$O(c^{-1/3})$ needed for the Proposition.
This means the dominant error in $|\alpha^{(c)} - \pi/2|$
comes from the BS increment $O(1/k^2)$,
not from the residual difference.
We record the sharp bound here for completeness;
it implies \emph{second-order discrete stability} of the
phase increment.
\end{remark}

\begin{proof}
By Lemma~C: $R_k(c) = E_k(c)$,
so $R_{k+1} - R_k = E_{k+1}(c) - E_k(c)$.

The Airy-connection error $E_k(c)$ is a $C^1$ function of the
Airy parameter $\mu_k := \chi_k/(2c^{4/3})$
(ORX \cite{OsipovRokhlinXiao}, Ch.~4, Eq.~(4.31)),
with $|\partial_\mu E| \le C_A c^{-1/3}$ uniformly
for $\mu_k \le C_A$.
The step in $\mu_k$:
\[
  |\mu_{k+1} - \mu_k|
  = \frac{|\chi_{k+1} - \chi_k|}{2c^{4/3}}
  = \frac{O(c^{1/3})}{2c^{4/3}}
  = O(c^{-1}).
\]
Hence:
\begin{equation}\label{eq:E-diff}
  |E_{k+1}(c) - E_k(c)|
  \le C_A c^{-1/3} \cdot O(c^{-1})
  = O(c^{-4/3}),
\end{equation}
giving $|R_{k+1} - R_k| = O(c^{-4/3})$.
\end{proof}

% ============================================================
\section{Proposition: Phase Increment}
\label{sec:prop}
% ============================================================

\begin{proposition}[Phase increment]
\label{prop:alpha}
For all $c \ge c_0(A)$:
\begin{equation}\label{eq:alpha-main}
  \sup_{1 \le k \le N-1}
  \left|\alpha^{(c)} - \frac{\pi}{2}\right|
  \le C_A\,c^{-1/3}.
\end{equation}
\end{proposition}

\begin{proof}
From Lemma~A:
$\alpha^{(c)} = \frac12(S^+(\chi_{k+1}) - S^+(\chi_k))
+ R_{k+1}(c) - R_k(c)$.
By Lemma~D: $|R_{k+1} - R_k| = O(c^{-4/3})$.

\textbf{Case 1: $k \ge K_0 := \lceil c^{1/6}\rceil$.}
By Paper~IV: $|\frac12(S^+(\chi_{k+1})-S^+(\chi_k)) - \pi/2|
= O(1/k^2) \le c^{-1/3}$ for $k \ge c^{1/6}$.
Hence $|\alpha^{(c)} - \pi/2| \le C c^{-1/3} + O(c^{-4/3})
= C_A' c^{-1/3}$.

\textbf{Case 2: $1 \le k < K_0$.}
For each fixed $k$, the BS half-increment converges to $\pi/2$
as $c \to \infty$ (ORX \cite{OsipovRokhlinXiao}, Ch.~4, Thm.~4.7),
with
$|\frac12(S^+(\chi_{k+1})-S^+(\chi_k)) - \pi/2|
\le C_k c^{-1/3}$
for $c \ge c_0(k)$.
Choosing $c_0(A)$ large enough to dominate all $C_k$
for $k \le K_0 = c^{1/6}$ (a finite set for each $c$),
the bound $|\alpha^{(c)} - \pi/2| \le C_A'' c^{-1/3}$
holds uniformly for all $1 \le k \le N-1$.
\end{proof}

% ============================================================
\section{Corollaries}
\label{sec:corollaries}
% ============================================================

\begin{corollary}[Non-degeneracy / Conjecture~6.1]
\label{cor:nondeg}
There exists $c_0 = c_0(A)$ such that for all
$c \ge c_0$ and $1 \le k \le N-1$:
\[
  \mathrm{dist}\bigl(\alpha^{(c)},\,\{0,\pi\}\bigr) \ge \frac{\pi}{4}.
\]
Conjecture~6.1 of Paper~VI \cite{paperVI} is proved.
\end{corollary}

\begin{proof}
$|\alpha^{(c)} - \pi/2| \le C_A c^{-1/3}$;
choose $c_0$: $C_A c_0^{-1/3} \le \pi/4$.
Then $\alpha^{(c)} \in [\pi/4, 3\pi/4]$.
\end{proof}

\begin{corollary}[H1 complete; Paper~VII unconditional]
\label{cor:paperVII}
Hypothesis~H1 of Paper~VII \cite{paperVII} is fully verified:
\begin{enumerate}[(a)]
  \item \textbf{H1 Summability:}
    Proposition~3.1(a) of Paper~VII (unconditional).
  \item \textbf{H1 Non-Degeneracy:}
    Corollary~\ref{cor:nondeg}.
\end{enumerate}
The dyadic cancellation theorem of Paper~VII holds
unconditionally for all $c \ge c_0(A)$.
Conjecture~6.1 of Paper~VI changes status from open to proved.
\end{corollary}

% ============================================================
\section{What Remains Open}
\label{sec:open}
% ============================================================

\begin{enumerate}[(O1)]
  \item \textbf{ass:gap:}
    $\lambda_l - \lambda_{l+1} \ge \kappa_0 c^{-1/3}$
    (next target).
  \item \textbf{B-strong:}
    $P_{kl} \le C_2 c^{1/2}$.
  \item \textbf{Amplitude regularity / Conjecture~6.2.}
  \item \textbf{Proposition~U' (Paper~VIII).}
\end{enumerate}

% ============================================================
\begin{thebibliography}{99}

\bibitem{paperIV}
  Anonymous,
  \textit{Semiclassical Equidistribution of PSWF Densities
  via Pr\"ufer Phase Analysis},
  Paper~IV, prolate-gram-coercivity, 2026.

\bibitem{paperVI}
  Anonymous,
  \textit{Galerkin Norm Estimates for the Prolate Phase Operator},
  Paper~VI, prolate-gram-coercivity, 2026.

\bibitem{paperVII}
  Anonymous,
  \textit{Dyadic Cancellation for PSWF Galerkin Operators},
  Paper~VII, prolate-gram-coercivity, 2026.

\bibitem{Olver}
  F.~W.~J.~Olver,
  \textit{Asymptotics and Special Functions},
  Academic Press, 1974.

\bibitem{OsipovRokhlinXiao}
  A.~Osipov, V.~Rokhlin, H.~Xiao,
  \textit{Prolate Spheroidal Wave Functions of Order Zero},
  Springer, 2013.

\end{thebibliography}

\end{document}
